\section{El modelo mejorará; las empresas de aplicaciones aportan el entorno y la experiencia}

La sección anterior trató el ecosistema de agentes; esta trata los límites de las empresas de aplicaciones. Ming Chaoping (明超平) dijo: «Confía siempre en que el modelo mejorará y en que el modelo no tiene nada que ver contigo». La frase no es pesimista; marca un límite para las empresas de aplicaciones: las empresas de modelos fundacionales crean personas más inteligentes, y las empresas de aplicaciones deben usar el entorno y la experiencia para que esos modelos más inteligentes se adapten a las necesidades reales de producción.

\begin{figure}[H]
\centering
\includegraphics[width=0.96\textwidth]{os-agent-living-system.png}
\caption{OS Agent: un sistema vivo. Los modelos fundacionales crean personas más inteligentes; las empresas de aplicaciones aportan el entorno y la experiencia. Diagrama conceptual propio, elaborado a partir de la conversación de 01:49:50--01:58:26.}
\end{figure}

\begin{knowledgebox}{Lectura de la figura: la empresa de aplicaciones no es una versión menor de la empresa de modelos}
En la figura, el Model se vuelve más inteligente, el Environment contiene los escenarios de uso, la Experience se adapta a las necesidades de producción, el OS Agent se convierte en un sistema vivo y el Feedback regresa de forma continua. El valor de la empresa de aplicaciones está en llevar la capacidad del modelo a flujos de trabajo repetibles y al entorno del usuario.
\end{knowledgebox}

\subsection{Exprimir la inteligencia con más eficiencia}

Ming Chaoping resume lo que observó en los últimos dos años como consumir tokens con más eficiencia y exprimir la inteligencia. En la era del Chatbot, el consumo de tokens dependía sobre todo del prompt del usuario; en la era del Agent, el sistema puede dividir las tareas por iniciativa propia, invocar herramientas y subagentes y generar productos intermedios, con lo que aumenta la velocidad de consumo de tokens. El problema es que ese consumo debe generar valor; de lo contrario, solo es costo.

\begin{importantbox}{La pregunta central para las empresas de aplicaciones}
No es «¿tengo un modelo?», sino «¿puedo convertir el modelo, con más eficiencia, en la capacidad de que el usuario complete sus tareas?». Esto exige que el producto entienda la tarea, el entorno, la retroalimentación, la distribución y el modelo de negocio, en lugar de solo envolver una API en otra capa.
\end{importantbox}

\subsection{Precios y modelo de negocio}

Antes se habló del entorno y del consumo de inteligencia; esta sección examina cómo se convierte en un modelo de negocio. En la entrevista se menciona que el precio de un producto Agent no tiene por qué basarse solo en el cobro por servicio. Pueden coexistir suscripciones, cobro por servicio, publicidad, incentivos para creadores y comisiones de plataforma. El juicio de Ming Chaoping es que el modelo de negocio debe generar beneficio mutuo para la plataforma y los creadores: que la empresa tenga utilidades suficientes para seguir desarrollando el producto y que los usuarios o creadores también reciban incentivos suficientes.

\begin{knowledgebox}{Términos clave: comercialización de aplicaciones de IA}
\noindent\begin{tabularx}{\textwidth}{p{0.20\textwidth}p{0.34\textwidth}X}
\toprule
Modelo & Escenario de uso & Riesgo \\
\midrule
Suscripción & Herramientas de uso frecuente, individuales o de equipo & Las empresas de modelos pueden incluirlas en sus paquetes y presionar el precio a la baja. \\
Cobro por servicio & Resultados de tarea definidos, como generar un sitio web, desplegarlo o ejecutar un flujo & Requiere controlar el costo y la calidad de la entrega. \\
Publicidad & Escenarios de contenido y distribución & Puede dañar la experiencia de creación. \\
Reparto con creadores & Ecosistema de la plataforma y comercio de obras & Requiere transacciones reales y gestión de derechos de autor. \\
Cobro a empresas & Colaboración en equipo, permisos, seguridad y despliegue & Ciclo de venta largo; el producto se vuelve pesado. \\
\bottomrule
\end{tabularx}
\end{knowledgebox}

\subsection{Por qué no basta con volverse to B}

Ming Chaoping mencionó que, si todas las empresas to C se convirtieran en empresas to B, habría resistencia generalizada. Detrás de esta frase está la cuestión de la ruta que sigue el emprendimiento en aplicaciones de IA. Muchos productos Agent se acercan a los servicios empresariales para tener ingresos seguros, pero si se vuelven completamente to B, el producto puede volverse pesado, el ciclo de venta se alarga, la comunidad de creación se debilita y se pierden los nuevos grupos de usuarios y los efectos de red.

Aquí está la tensión de YouWare: necesita comercializarse, pero no puede encerrarse demasiado pronto en el software empresarial tradicional. La creación con coding, la publicación de obras, los Vibe Coder y la red de Agent se parecen más a un ecosistema to C o prosumer; los permisos, la identidad, la colaboración en equipo y los pagos, en cambio, la empujan hacia el lado B. Una estrategia de producto exigente debería conservar las posibilidades de ambos lados, en lugar de sacrificar los accesos futuros a cambio de ingresos de corto plazo.

\begin{importantbox}{Elección de ruta: primero ser el entorno, después ser el producto que se vende}
Si YouWare se convierte primero en un entorno de creación, tendrá la oportunidad de acumular obras, usuarios, invocaciones de Agent y efectos de red; si demasiado pronto se dedica solo a entregas empresariales, puede obtener ingresos, pero le será difícil construir AgentRank y posicionarse en la mente de nuevos grupos de usuarios.
\end{importantbox}

\subsection{Resumen de la sección}

El valor de las empresas de aplicaciones no está en copiar los modelos fundacionales, sino en construir el entorno, la experiencia y la retroalimentación. La pregunta de largo plazo para YouWare es si puede pasar de ser un punto de entrada a la creación con coding a ser un sistema vivo que consuma inteligencia de forma continua, produzca obras y genere efectos de red.
